% !TeX program = pdflatex
% !TeX encoding = UTF-8
% Preamble.
% !TeX program = pdflatex
% !TeX encoding = UTF-8
% Preamble.
\documentclass[12pt]{article}

% Packages.
\usepackage[utf8]{inputenc}     % Unicode support (Umlauts etc.)
\usepackage{amsmath, amssymb}   % Advanced math typesetting
\usepackage{stmaryrd, bbm}
\usepackage[french]{babel}      % Change hyphenation rules
\usepackage[T1]{fontenc}
\usepackage{geometry}
\usepackage{listings}           % Source code formatting and highlighting
\usepackage{csquotes}
\usepackage{interval}
\usepackage{pifont}

% Pages.
\sloppy 
\geometry{a4paper, total={17cm, 25cm}}
\renewcommand{\baselinestretch}{1.4}

%% Real set.
\newcommand{\R}{\mathbb{R}}
%% Length of vector.
\newcommand{\norm}[1]{\left\lVert\ #1\;\right\rVert}
%% Parenthesis.
\newcommand{\of}[1]{\left(#1\right)}

% Data.
\title{MP2I - Sommes, Produits et systèmes linéaires.}
\author{Hector GARLATTI.}
\date{\today{}}

\begin{document}
\maketitle

$\bigstar$ Fiche de synthèse.

\section{Notations.}

Soit des familles $(u_i)_{i \ge 0}, (v_i)_{i \ge 0}$.

\begin{itemize}
    \item Changement d'indice : soit $k, i \in \mathbb{Z}, \; i = k + c$:
$$
\sum_{k = a}^{b} u_k = \sum_{i = a + c}^{b + c} u_{i - c}
$$
    \item Linéarité de la somme : $$
\sum_{k=a}^{b} (\lambda u_k + v_k) = 
\lambda \sum_{k=a}^{b} u_k + \sum_{k=a}^{b} v_k
$$
\end{itemize}

\section{Propriétés.}

Soit des familles $(u_i)_{i \ge 0}, (v_i)_{i \ge 0}, 
\alpha \in \mathbb{R}, I \subset \mathbb{R}$  :

\begin{itemize}
    \item Somme repétées d'un scalaire : $$
\sum_{i \in I} \alpha = \alpha \vert I \vert
$$
    \item Relation de Chasles : si $p \le q \le r$: $$
\sum_{k=p}^{r} u_k = \sum_{k=p}^{q} u_k + \sum_{k=q+1}^{r} u_k   
$$
    \item Sommation par paquets : si $n \in \mathbb{N}, 
I = I_1 \cup I_2 \cup \dots \cup I_n$ tous disjoints : $$
\sum_{i \in I} u_i = \sum_{k=1}^{n} \sum_{i \in I_k} u_i
$$
    \item Somme \enquote{téléscopiques} : $$
\sum_{k=p}^{q} (u_{k+1} - u_k) = u_{q + 1} - u_p  
$$
    \item Double somme et produit : si $w_{i,j} = u_i v_i$, $$
\sum_{i=1}^{n} \sum_{j=1}^{n} w_{i, j} = 
\left(\sum_{i=1}^{n} u_i\right) \left(\sum_{j=1}^{n} v_i \right)     
$$
\end{itemize}

\section{Formules remarquables.}

Soit $q \in \mathbb{C}$.

\begin{itemize}
    \item Somme des termes d'une suite géométrique : $$
\sum_{k = a}^{b} = q^a \frac{1 - q^{b - a + 1}}{1 - q}  
$$
    \item Somme des $n \in \mathbb{N}$ premiers entier : $$
\sum_{k = 0}^{n} = \frac{n(n + 1)}{2}  
$$ 
    \item Somme des $n \in \mathbb{N}$ premiers carrés : $$
\sum_{k = 0}^{n} = \frac{n(n + 1)(2n + 1)}{6}  
$$
    \item $3^{\text{ème}}$ identité remarquable généralisée :
si $a, b \in \mathbb{R}, n \in \mathbb{N}$ : $$
a^n - b^n = (a - b)\sum_{k = 1}^{n} a^{n - k}b^{k - 1}
$$
    \item Carré d'une somme : $$
\left(\sum_{i=1}^{n} u_i \right)^2 =
\sum_{i=1}^{n} u_i^2 + 2\sum_{1 \le i < j \le n} u_i u_j  
$$
    \item Inégalié de Cauchy-Schwarz : $$
\left(\sum_{i=1}^{n} u_i v_i \right)^2 =
\left(\sum_{i=1}^{n} u_i^2 \right) 
\left(\sum_{i=1}^{n} v_i^2 \right)  
$$
\end{itemize}

\section{Coefficient binomiaux.}

Soit $n, k \in \mathbb{N}, k > n$ :
\begin{itemize}
    \item Coefficient bionmial : $$
\begin{pmatrix} n \\ k \end{pmatrix} = 
\frac{n!}{k!(n - k)!} 
$$
    \item Formule du binome : $$
(a + b)^n 
= \sum_{k=0}^{n} \begin{pmatrix} n \\ k \end{pmatrix} a^k b^{n-k}
$$
\end{itemize}

\section{Systèmes linéaires.} 

Pour résoudre un système à $n$ équations et $p$ inconnues,
il faut se rammener à un système diagonal à 
coefficients  diagonaux non-nuls en utilisant
les operations élémentaire qui préservent l'équivalence :
\begin{itemize}
    \item Multiplication : Soit $\lambda \in \mathbb{R}^{*}$, $$
    L_i \leftarrow \lambda L_i
$$
    \item Échanger : $$ L_i \leftrightarrow L_j$$
    \item Ajouter : $$ L_i \leftarrow L_i + L_j$$
\end{itemize}

\end{document}
